% !TeX root = ../main.tex

\section{การทดลองที่ 3: ผลกระทบของ $Q$ และอัตราส่วน $Q/R$
         ต่อพฤติกรรมของ Scalar KF}
% \textit{ครอบคลุม Criteria: Part 2 --- ออกแบบการทดลองแสดงผลของการปรับ $Q$,
% อธิบาย $Q/R$ Trade-off, เลือก Visualization ที่เหมาะสม (รวม 2.5 คะแนน)}

%% ─────────────────────────────────────────────────────────────────────────────
\subsection{บทนำ}

จากการทดลองที่ 2 ได้ Scalar Kalman Filter ที่ทำงานได้แล้ว พร้อมค่า
Measurement Noise Covariance $R = 0.0337\ \text{cm}^2$ ที่คำนวณได้จาก Variance
ของ Raw Data จริง อย่างไรก็ตาม การเลือกค่า $Q$ เป็นหนึ่งในการตัดสินใจทางวิศวกรรม
ที่สำคัญที่สุดในการออกแบบ Kalman Filter เนื่องจาก $Q$ ควบคุมสมดุลระหว่างสองเป้าหมาย
ที่ขัดแย้งกัน ได้แก่ \textbf{Smoothness} (ความเรียบของสัญญาณออก)
และ \textbf{Responsiveness} (ความไวในการติดตามการเปลี่ยนแปลง)

กลไกหลักที่ $Q$ ส่งผลผ่านคือสมการ Covariance Prediction:
\begin{equation}
    P_k^{(-)} = P_{k-1} + Q
    \label{eq:cov_predict_exp3}
\end{equation}
ทุกรอบที่ระบบทำงาน $Q$ จะ``ฉีด''ความไม่แน่นอนเข้าไปใน $P_k^{(-)}$
ทำให้ Kalman Gain $K_k$ มีค่าสูงขึ้น และ Filter หันมาเชื่อถือ Measurement มากขึ้น
ในทางกลับกัน $Q$ ที่ต่ำมากจะทำให้ $P_k^{(-)}$ ลู่เข้าศูนย์
ส่งผลให้ Filter หยุดรับข้อมูลใหม่จาก Sensor

การเลือกทดสอบ $Q \in \{1\times10^{-3},\ 1\times10^{-5},\ 1\times10^{-6}\}$
ตามแนวทางของ Al~Tahtawi (2018) เนื่องจากช่วงนี้ครอบคลุมอัตราส่วน $Q/R$
ตั้งแต่ระดับที่ Filter เชื่อ Measurement มากกว่า Model ($Q \gg R$)
ไปจนถึงระดับที่ Filter เชื่อ Model มากกว่า Measurement ($Q \ll R$)
ทำให้สามารถสังเกตพฤติกรรมสุดโต่งทั้งสองด้านได้ในการทดลองเดียวกัน
โดยใช้ $Q = 1\times10^{-5}$ เป็น Nominal Value เนื่องจากมีขนาดใกล้เคียงกับ
Variance ของ Raw Data ที่ระยะใกล้ ($s^2 = 0.0050\ \text{cm}^2$ ที่ระยะ 10 cm)
ซึ่งสะท้อนระดับความไม่แน่นอนของ Model ที่สมเหตุสมผลสำหรับ Sensor นี้

งานวิจัย Al~Tahtawi (2018) ได้ทดสอบค่า $Q$ สามระดับดังกล่าวบน HC-SR04
และพบว่า Output Noise Variance มีความสัมพันธ์แบบ Exponential กับขนาดของ $Q$
การทดลองนี้จึงทำซ้ำการทดสอบดังกล่าวบนแพลตฟอร์ม STM32 Nucleo G474RE
เพื่อยืนยันผลและเลือกค่า $Q$ ที่เหมาะสมที่สุดสำหรับการนำไปใช้จริง

%% ─────────────────────────────────────────────────────────────────────────────
\subsection{วัตถุประสงค์}

\begin{itemize}
    \item เพื่อศึกษาว่าการเปลี่ยนแปลงค่า $Q$ ส่งผลต่อพฤติกรรมของ KF อย่างไร
          ทั้งในสถานะคงที่ (Steady State) และเมื่อมี Step Change
    \item เพื่อทำความเข้าใจ Trade-off ระหว่าง Smoothness และ Responsiveness
          ที่ถูกควบคุมโดยอัตราส่วน $Q/R$
    \item เพื่อวัดผล Quantitative ของ Smoothness (Output Variance)
          และ Responsiveness (Response Time) สำหรับแต่ละ $Q$
    \item เพื่อเลือกและ Justify ค่า $Q$ ที่เหมาะสมที่สุดสำหรับ
          ระบบ HC-SR04 บน STM32 Nucleo G474RE
\end{itemize}

%% ─────────────────────────────────────────────────────────────────────────────
\subsection{สมมติฐาน}

\begin{itemize}
    \item $Q$ ที่มากขึ้นจะทำให้ Filter ตอบสนองเร็วขึ้น
          (ติดตาม Step Change ได้ไวขึ้น) แต่สัญญาณออกจะมี Noise มากขึ้น
    \item $Q$ ที่น้อยลงจะให้สัญญาณออกเรียบขึ้น
          แต่เกิด Lag Error เมื่อระยะทางจริงเปลี่ยนอย่างกะทันหัน
    \item Output Noise Variance จะมีความสัมพันธ์เป็นเชิงบวกกับขนาดของ $Q$
          ตามที่ Al Tahtawi (2018) รายงานไว้
    \item มีช่วงค่า $Q$ ที่ให้ผลสมดุลระหว่าง Smoothness และ Responsiveness
          สำหรับ Sensor และ Application นี้โดยเฉพาะ
\end{itemize}

%% ─────────────────────────────────────────────────────────────────────────────
\subsection{ตัวแปร}

\begin{itemize}
    \item \textbf{ตัวแปรอิสระ:} ค่า $Q$ ที่ทดสอบ 3 ระดับ
          ตาม Al Tahtawi (2018) ได้แก่
          $Q = 1\times10^{-3}$, $Q = 1\times10^{-5}$, และ $Q = 1\times10^{-6}$
    \item \textbf{ตัวแปรตาม:} Variance ของ KF Output (วัด Smoothness),
          Response Time หลัง Step Change (วัด Responsiveness),
          Steady-State Kalman Gain $K_{k,\mathrm{ss}}$
    \item \textbf{ตัวแปรควบคุม:} ค่า $R = 0.0337$~cm$^2$ คงที่
          (จากการทดลองที่ 2), Sampling Time $T_s = 60$~ms คงที่,
          Scenario การทดสอบเหมือนกันทุกค่า $Q$
\end{itemize}

%% ─────────────────────────────────────────────────────────────────────────────
\subsection{เอกสารและงานวิจัยที่เกี่ยวข้อง}

\subsubsection{ความหมายทางกายภาพของ Process Noise Covariance $Q$}

$Q$ คือพารามิเตอร์ที่สะท้อน \textbf{ความไม่เชื่อมั่นใน Dynamic Model} ของระบบ
หากระบบจริงมีพฤติกรรมที่ Model ไม่ได้จำลอง (เช่น การเคลื่อนที่ของเป้าหมายในระบบ Static
หรือแรงรบกวนที่ Model ไม่ครอบคลุม) $Q$ ควรมีค่าสูงขึ้นเพื่อให้ Filter
มีความยืดหยุ่นพอที่จะปรับตัวตามสถานการณ์จริง ในทางตรงข้าม
หากระบบจริงสอดคล้องกับ Model อย่างแม่นยำ $Q$ ควรมีค่าต่ำ

แนวคิดนี้แตกต่างจาก $R$ ซึ่งสะท้อนความไม่เชื่อมั่นใน \textbf{Sensor}
ความสัมพันธ์ระหว่าง $Q$ และ $R$ จึงกำหนดว่า Filter จะ ``เชื่อ Model''
หรือ ``เชื่อ Sensor'' มากกว่ากัน ซึ่งควบคุมโดยอัตราส่วน $Q/R$

\subsubsection{กลไกผ่านสมการ Covariance Prediction}

จากสมการ Covariance Prediction ของ Scalar KF
\begin{equation}
    P_k^{(-)} = P_{k-1} + Q
    \label{eq:cov_pred_detail}
\end{equation}

$Q$ บวกเข้าไปใน $P_k^{(-)}$ ทุกรอบ ทำให้ $P_k^{(-)}$ ไม่สามารถลู่เข้าศูนย์ได้
ผลที่ได้คือ Kalman Gain
\begin{equation}
    K_k = \frac{P_k^{(-)}}{P_k^{(-)} + R}
    \label{eq:scalar_gain_exp3}
\end{equation}
จะ Converge ไปสู่ค่า Steady-State $K_{k,\mathrm{ss}}$ ที่ไม่เป็นศูนย์
ค่า $K_{k,\mathrm{ss}}$ หาได้จากการแก้สมการกำลังสองของ $P_\mathrm{ss}$
\begin{equation}
    P_\mathrm{ss}^2 + Q \cdot P_\mathrm{ss} - Q \cdot R = 0
    \label{eq:p_ss_quadratic}
\end{equation}
ซึ่งมีรากบวกเป็น
\begin{equation}
    P_\mathrm{ss} = \frac{-Q + \sqrt{Q^2 + 4QR}}{2}
    \label{eq:p_ss_solution}
\end{equation}

ที่ค่า $R = 0.0337$~cm$^2$ ที่ได้จากการทดลองที่ 2
$K_{k,\mathrm{ss}}$ ของ $Q$ ทั้งสามค่ามีดังต่อไปนี้

\begin{table}[htbp]
    \centering
    \caption{ค่า Steady-State Kalman Gain สำหรับแต่ละ $Q$
             (คำนวณจากสมการ~\eqref{eq:p_ss_solution}, $R = 0.0337$~cm$^2$)}
    \label{tab:kss_theory}
    \begin{tabular}{cccc}
        \hline
        \textbf{ค่า $Q$} & \textbf{$Q/R$} & \textbf{$P_\mathrm{ss}$ (cm$^2$)} & \textbf{$K_{k,\mathrm{ss}}$} \\
        \hline
        $1\times10^{-3}$ & $2.97\times10^{-2}$ & $0.00533$ & $0.136$ \\
        $1\times10^{-5}$ & $2.97\times10^{-4}$ & $0.00058$ & $0.017$ \\
        $1\times10^{-6}$ & $2.97\times10^{-5}$ & $0.00018$ & $0.005$ \\
        \hline
    \end{tabular}
\end{table}

จากตารางที่~\ref{tab:kss_theory} จะเห็นได้ว่าทั้งสาม $Q$ มีอัตราส่วน $Q/R \ll 1$
เนื่องจาก $R = 0.0337$~cm$^2$ มีขนาดใหญ่กว่า $Q$ ทุกค่ามาก
การจัดอันดับของ $K_{k,\mathrm{ss}}$ และพฤติกรรมที่แตกต่างกันจึงมาจาก
\textbf{ขนาดสัมพัทธ์ระหว่าง $Q$ ด้วยกันเอง} ไม่ใช่จากความสัมพันธ์
$Q$ vs.\ $R$ โดยตรง

%% ─────────────────────────────────────────
\subsubsection{Trade-off ระหว่าง Smoothness และ Responsiveness}

หัวใจสำคัญของการออกแบบ Kalman Filter คือการหาจุดสมดุล (Trade-off) ระหว่างการกำจัดสัญญาณรบกวนให้ได้มากที่สุด (Smoothness) และการรักษาความสามารถในการติดตามการเปลี่ยนแปลงของสัญญาณจริงให้ได้เร็วที่สุด (Responsiveness) พฤติกรรมนี้ถูกควบคุมโดยตรงจากอัตราส่วน $Q/R$ ซึ่งเป็นตัวกำหนดค่าอัตราขยายคาลมานในสภาวะคงตัว (Steady-State Kalman Gain, $K_{k,ss}$) 

ตารางที่ \ref{tab:q_tradeoff} สรุปความสัมพันธ์ระหว่างพารามิเตอร์ที่ใช้ทดสอบกับผลลัพธ์เชิงคุณภาพที่สังเกตได้จากกราฟการทดลอง

\begin{table}[htbp]
    \centering
    \caption{สรุป Trade-off ของผลการตอบสนองในแต่ละระดับค่า $Q$}
    \label{tab:q_tradeoff}
    \begin{tabular}{ccccc}
        \hline
        \textbf{ค่า $Q$} & \textbf{อัตราส่วน $Q/R$} & \textbf{$K_{k,ss}$} & \textbf{Smoothness} & \textbf{Response Time} \\
        \hline
        $1 \times 10^{-3}$ & $2.97 \times 10^{-2}$ & สูง ($\approx 0.136$) & ต่ำ  & เร็ว \\
        $1 \times 10^{-5}$ & $2.97 \times 10^{-4}$ & กลาง ($\approx 0.017$) & กลาง & กลาง \\
        $1 \times 10^{-6}$ & $2.97 \times 10^{-5}$ & ต่ำ ($\approx 0.005$)  & สูง  & ช้า \\
        \hline
    \end{tabular}
\end{table}



จากตารางที่ \ref{tab:q_tradeoff} สามารถอภิปรายพฤติกรรมของระบบได้ดังนี้:

\begin{itemize}
    \item \textbf{กรณี $Q = 1 \times 10^{-3}$ (เน้นการตอบสนองเร็ว):} เมื่อตั้งค่า $Q$ ให้มีขนาดใหญ่เมื่อเทียบกับ $R$ ระบบจะตีความว่าความคลาดเคลื่อนของโมเดลมีสูง จึงให้น้ำหนักความเชื่อมั่นไปที่ค่าจากการวัด (Measurement) มากกว่า สังเกตได้จากค่า $K_{k,ss}$ ที่มีค่าสูง ทำให้ระบบตอบสนองต่อการเปลี่ยนระยะ (Step Change) ได้ทันที แต่ข้อเสียคือสัญญาณที่กรองได้จะยังมีสัญญาณรบกวนปะปนอยู่มาก (Smoothness ต่ำ)
    
    \item \textbf{กรณี $Q = 1 \times 10^{-6}$ (เน้นความเรียบเนียน):} เมื่อตั้งค่า $Q$ ให้มีขนาดเล็กมาก ระบบจะเชื่อมั่นในโมเดลทางคณิตศาสตร์ (ซึ่งพยากรณ์ว่าระยะทางคงที่) มากกว่าเซนเซอร์ ค่า $K_{k,ss}$ จึงมีค่าต่ำมากเพียง $0.005$ ผลลัพธ์คือสัญญาณที่ได้จะมีความเรียบเนียนสูงมาก ทรงตัวเป็นเส้นตรง แต่เมื่อเกิดการเคลื่อนที่จริง ระบบจะเกิดความล่าช้า (Lag/Delay) ในการปรับตัวเข้าหาค่าใหม่ (Response Time ช้า)
    
    \item \textbf{กรณี $Q = 1 \times 10^{-5}$ (จุดสมดุล):} เป็นค่า Nominal ที่ให้ความสมดุลระหว่างสองขั้ว ระบบสามารถกรอง Measurement Noise ออกได้ในระดับที่น่าพอใจ ทำให้สัญญาณนิ่งขึ้นอย่างเห็นได้ชัด ในขณะเดียวกันก็ยังมีค่า $K_{k,ss}$ ที่สูงเพียงพอที่จะติดตามการเปลี่ยนแปลงระยะทางได้อย่างกระฉับกระเฉง โดยไม่ก่อให้เกิด Phase Lag ที่สังเกตเห็นได้ชัดเจนใน Application นี้
\end{itemize}

การวิเคราะห์นี้ยืนยันว่าการจูนพารามิเตอร์ $Q$ ไม่ใช่การหาค่าที่ "ถูกต้องที่สุด" แต่เป็นการพิจารณาหาค่าที่เหมาะสมกับข้อกำหนดการใช้งาน (Application Requirements) ว่าต้องการกรองสัญญาณรบกวนมากน้อยเพียงใด และยอมรับความล่าช้าของระบบได้มากแค่ไหน

\subsubsection{สภาวะ Overconfident Filter ($Q \to 0$)}

หาก $Q = 0$ สมการ Covariance Prediction กลายเป็น $P_k^{(-)} = P_{k-1}$
ซึ่งทำให้ $P_k$ ลู่เข้าศูนย์เมื่อมีการวัดซ้ำๆ ส่งผลให้ $K_k \rightarrow 0$
Filter จะหยุดรับข้อมูลใหม่จาก Sensor ทั้งหมด และยึดติดกับ Prediction ของ Model
เพียงอย่างเดียว สภาวะนี้เรียกว่า \textbf{Overconfident Filter}
ซึ่งเป็นอันตรายในระบบ Dynamic ที่สถานะจริงเปลี่ยนแปลงอยู่ตลอดเวลา

\subsubsection{งานวิจัย Al Tahtawi (2018) ที่เกี่ยวข้องโดยตรง}

Al Tahtawi (2018) ทดสอบ $Q$ สามระดับ ได้แก่ $Q = 1\times10^{-3}$,
$1\times10^{-5}$, $1\times10^{-6}$ บน ATMega328p โดยกำหนด $R = 2.92\times10^{-3}$
ผลสำคัญจาก Figure 6--10 มีดังนี้

\begin{itemize}
    \item \textbf{Figure 6 (Steady State):} $Q = 1\times10^{-6}$ ให้ Output
          เรียบที่สุดที่สถานะคงตัว ขณะที่ $Q = 1\times10^{-3}$ ให้ Output
          ใกล้เคียง Raw Data (a)
    \item \textbf{Figure 8--9 (Step Change):} $Q = 1\times10^{-3}$ ติดตาม
          Step Change ได้เร็วที่สุด ขณะที่ $Q = 1\times10^{-6}$
          มี Response Time ช้าที่สุด (b-c)
    \item \textbf{Figure 10 (Variance Analysis):} Output Noise Variance
          มีความสัมพันธ์เป็นเชิงบวกกับขนาดของ $Q$ โดย
          $Q = 1\times10^{-3}$ ให้ Variance สูงสุดและ $Q = 1\times10^{-6}$
          ให้ Variance ต่ำสุด (d)
\end{itemize}

\begin{figure}[H]
    \centering
    % แถวที่ 1
    \begin{minipage}{0.48\textwidth}
        \centering
        \includegraphics[width=\textwidth]{/Users/thadzy/Documents/01_Active/University/FRA233-Control/Lab3/FRA233_Lab3_Report/images/fig6.png}
        \vspace{0.2cm}
        \centerline{\small (a) Figure 6: Steady State Response}
    \end{minipage}\hfill
    \begin{minipage}{0.48\textwidth}
        \centering
        \includegraphics[width=\textwidth]{/Users/thadzy/Documents/01_Active/University/FRA233-Control/Lab3/FRA233_Lab3_Report/images/fig8.png}
        \vspace{0.2cm}
        \centerline{\small (b) Figure 8: Step Change Tracking}
    \end{minipage}
    
    \vspace{0.5cm} % เว้นระยะห่างระหว่างแถว
    
    % แถวที่ 2
    \begin{minipage}{0.48\textwidth}
        \centering
        \includegraphics[width=\textwidth]{/Users/thadzy/Documents/01_Active/University/FRA233-Control/Lab3/FRA233_Lab3_Report/images/fig9.png}
        \vspace{0.2cm}
        \centerline{\small (c) Figure 9: Step Change with Offset}
    \end{minipage}\hfill
    \begin{minipage}{0.48\textwidth}
        \centering
        \includegraphics[width=\textwidth]{/Users/thadzy/Documents/01_Active/University/FRA233-Control/Lab3/FRA233_Lab3_Report/images/fig10.png}
        \vspace{0.2cm}
        \centerline{\small (d) Figure 10: Variance Analysis}
    \end{minipage}
    
    \caption{ผลการตอบสนองของ Kalman Filter ต่อค่า Process Noise Covariance ($Q$) ในสภาวะต่างๆ (อ้างอิงจาก Al Tahtawi, 2018)}
    \label{fig:altahtawi_results}
\end{figure}

ข้อแตกต่างจากการทดลองนี้คือ Al Tahtawi (2018) ใช้ $R = 2.92\times10^{-3}$
ซึ่งเล็กกว่า $Q = 1\times10^{-3}$ ทำให้ $Q/R > 1$ ในกรณีนั้น
แต่ในการทดลองของเราใช้ $R = 0.0337$~cm$^2$ ซึ่งใหญ่กว่า $Q$ ทุกค่า
ดังนั้น $Q/R < 1$ ในทุกกรณี อย่างไรก็ตาม \textbf{ทิศทางของ Effect ยังเหมือนกัน}
คือ $Q$ ที่ใหญ่กว่าทำให้ Filter ตอบสนองเร็วกว่าเสมอ

%% ─────────────────────────────────────────────────────────────────────────────
\subsection{ขั้นตอนการทดลอง}

การทดลองนี้ดำเนินการใน 2 Scenario หลัก โดยใช้ KF Library จากการทดลองที่ 2
เปลี่ยนเฉพาะค่า $Q_0$ ระหว่าง Run โดยคง $R = 0.0337$~cm$^2$ และ $T_s = 60$~ms
ตามเดิม

\textbf{Scenario A: ระยะคงที่ (Steady-State Analysis)}

\begin{enumerate}
    \item วางวัตถุแบนที่ระยะ 10~cm จาก Sensor ตั้งฉากกับเซนเซอร์
    \item รัน KF ด้วย $Q$ ทั้ง 3 ค่า แยก Run กัน บันทึก 10 วินาทีต่อ Run
    \item Log ค่า Raw $z_k$, KF Output $\hat{x}_k$, Kalman Gain $K_k$,
          และ Error Covariance $P_k$ ทุก Timestamp
    \item คำนวณ Output Variance ของแต่ละ Run ด้วยสมการ Sample Variance
          ตามที่นิยามในการทดลองที่~2 โดยใช้สูตร
          $s^2 = \sum_{i=1}^{N}(z_i - \bar{z})^2 / (N-1)$
\end{enumerate}

\textbf{Scenario B: MSD Free Vibration (Responsiveness Analysis)}
 
การทดลอง Scenario B ใช้ชุด Mass-Spring-Damper จริงที่ติดตั้ง HC-SR04
ไว้ด้านล่างชี้ขึ้นเพื่อวัดระยะของมวลในแนวดิ่ง โดยมีขั้นตอนดังนี้
 
\begin{enumerate}
    \item ปล่อยให้มวลอยู่ที่ตำแหน่ง Equilibrium ($x = 0$)
          ซึ่งกำหนดเป็นตำแหน่งที่ แรงสปริง $Kx$ พอดีสมดุลกับน้ำหนักมวล $mg$
          ยืนยันว่า Sensor อ่านค่า Baseline คงที่ก่อนเริ่มทดลอง
    \item ดึงมวลออกจาก Sensor ในแนวดิ่ง (ดึงขึ้น) เป็นระยะ
          $x_0 = +7.2$~cm จากจุด Equilibrium
          แล้วปล่อยโดยไม่ออกแรง ($\dot{x}(0) = 0$)
          ทำให้ระบบเกิด Free Vibration
    \item บันทึกข้อมูลตั้งแต่ช่วงที่ปล่อยมวลจนกระทั่ง Oscillation
          สงบลง (อย่างน้อย 10~วินาที) ส่งข้อมูลผ่าน Serial 115200~baud
          ในรูปแบบ \texttt{t\_ms, raw, filt\_Q1e-3, filt\_Q1e-5, filt\_Q1e-6}
          โดย KF ทั้ง 3 ค่า $Q$ รันพร้อมกันในลูปเดียว
    \item ทำซ้ำ 3 Trials โดยให้มวลกลับสู่ Equilibrium ก่อนเริ่ม Trial ถัดไป
\end{enumerate}
 
\textbf{การคำนวณ Settling Time สำหรับ Free Vibration}
 
เนื่องจาก Scenario B เป็น Free Vibration ไม่ใช่ Step Change
จึงไม่สามารถใช้นิยาม Response Time แบบเดิม (เข้าใกล้ Setpoint ใหม่ภายใน 5\%)
ได้โดยตรง การทดลองนี้จึงนิยาม \textbf{Settling Time} ดังนี้
 
\begin{itemize}
    \item คำนวณ \textbf{Initial Amplitude} $A_0$ จากค่า Maximum Deviation
          ของ KF Output จาก Equilibrium ในช่วง 30\% แรกของการบันทึก
          \begin{equation}
              A_0 = \max_{k \,\in\, [1,\, 0.3N]} \left|\hat{x}_k - \bar{x}_{\mathrm{eq}}\right|
              \label{eq:init_amp}
          \end{equation}
          โดยที่ $\bar{x}_{\mathrm{eq}}$ คือค่าเฉลี่ย KF Output ใน 20\% สุดท้าย
          ซึ่งระบบน่าจะนิ่งแล้ว
    \item กำหนด \textbf{Settling Threshold} เป็น 5\% ของ $A_0$
          \begin{equation}
              \delta = 0.05 \times A_0
              \label{eq:settle_threshold}
          \end{equation}
    \item \textbf{Settling Time} $t_s$ คือเวลาที่ KF Output เข้าสู่แถบ
          $[\bar{x}_{\mathrm{eq}} - \delta,\; \bar{x}_{\mathrm{eq}} + \delta]$
          และอยู่ในแถบนั้นตลอดจนสิ้นสุดการบันทึก
          \begin{equation}
              t_s = t_{k^*} - t_1
              \quad \text{โดยที่} \quad
              k^* = \min\left\{k : \left|\hat{x}_j - \bar{x}_{\mathrm{eq}}\right| \leq \delta
              \;\;\forall j \geq k \right\}
              \label{eq:settling_time}
          \end{equation}
    \item $Q$ ที่มี $t_s$ \textbf{น้อยกว่า} หมายความว่า Filter ปรับตัวตาม
          พลศาสตร์ของระบบได้เร็วกว่า จึงวัด Responsiveness ได้ในบริบท Dynamic
\end{itemize}
 
\begin{enumerate}
    \setcounter{enumi}{4}
    \item วิเคราะห์ผลโดยคำนวณ Settling Time ของ KF Output
          แต่ละค่า $Q$ ตามนิยามในสมการ~\eqref{eq:settling_time}
          ด้วย MATLAB โดยใช้ Envelope Method ข้างต้น
    \item Plot Raw $z_k$ และ KF Output $\hat{x}_k$ ของ $Q$ ทั้ง 3 ค่า
          บนกราฟเดียวกัน เพื่อเปรียบเทียบ Smoothness และ Tracking
          ในสภาวะ Dynamic
\end{enumerate}

%% ─────────────────────────────────────────────────────────────────────────────
%% ส่วนที่เพิ่ม/แทนที่ใน sections/04_exp3.tex
%% เฉพาะส่วน ผลการทดลอง, สรุป, อภิปราย
%% (ส่วน บทนำ, วัตถุประสงค์, สมมติฐาน, ตัวแปร, เอกสาร, ขั้นตอน มีอยู่แล้ว)
%% ─────────────────────────────────────────────────────────────────────────────

\subsection{ผลการทดลอง}

\subsubsection{Scenario A: ระยะคงที่ --- Variance Reduction}

ผลการคำนวณ Output Variance ของ KF แต่ละค่า $Q$ เปรียบเทียบกับ Raw Data
จากการทดลอง 3 Trials แสดงในตารางที่~\ref{tab:exp3_scA_variance}

\begin{table}[htbp]
    \centering
    \caption{Variance ของ Raw Data และ KF Output ที่ระยะคงที่ 10~cm
             (Scenario A, 3 Trials, 152 Samples ต่อ Run)}
    \label{tab:exp3_scA_variance}
    \begin{tabular}{clccr}
        \hline
        \textbf{Trial}
            & \textbf{ค่า $Q$}
            & \textbf{Var Raw (cm$^2$)}
            & \textbf{Var KF (cm$^2$)}
            & \textbf{Noise ลดลง (\%)} \\
        \hline
        \multirow{3}{*}{1}
            & $1\times10^{-3}$ & 0.1712 & 0.0290 & 83.0 \\
            & $1\times10^{-5}$ & 0.1732 & 0.0287 & 83.4 \\
            & $1\times10^{-6}$ & 0.2266 & 0.0430 & 81.0 \\
        \hline
        \multirow{3}{*}{2}
            & $1\times10^{-3}$ & 0.1487 & 0.0374 & 74.8 \\
            & $1\times10^{-5}$ & 0.2158 & 0.0287 & 86.7 \\
            & $1\times10^{-6}$ & 0.2100 & 0.0196 & 90.7 \\
        \hline
        \multirow{3}{*}{3}
            & $1\times10^{-3}$ & 0.0765 & 0.0217 & 71.7 \\
            & $1\times10^{-5}$ & 0.0001 & 0.0000 & 82.6 \\
            & $1\times10^{-6}$ & 0.0119 & 0.0001 & 98.9 \\
        \hline
        \multirow{3}{*}{\textbf{Mean}}
            & $1\times10^{-3}$ & 0.1321 & 0.0294 & \textbf{76.5} \\
            & $1\times10^{-5}$ & 0.1297 & 0.0191 & \textbf{84.2} \\
            & $1\times10^{-6}$ & 0.1495 & 0.0209 & \textbf{90.2} \\
        \hline
    \end{tabular}
\end{table}
 
\noindent\textit{หมายเหตุ:} Trial 3 ที่ $Q = 1\times10^{-5}$
มีค่า Raw Variance ต่ำผิดปกติ ($4.9\times10^{-5}$~cm$^2$)
เนื่องจาก Sensor อ่านค่าได้นิ่งมากในช่วงนั้น โดยสัญญาณแกว่งอยู่ในช่วง
เพียง $\pm 0.015$~cm ซึ่งอยู่ในระดับ Quantization ของ HC-SR04
จึงรายงานเป็น 0.0001~cm$^2$ ตามค่าจริงที่คำนวณได้ โดยไม่ได้ตัดทิ้ง

\begin{figure}[htbp]
    \centering
    % TODO: แทนด้วยกราฟจริงจาก MATLAB
    \includegraphics[width=0.85\textwidth]{exp3_scenarioA.png}
    \caption{Time-Domain Plot (Scenario A): Raw Data $z_k$ และ KF Output $\hat{x}_k$
             สำหรับ $Q \in \{1\times10^{-3},\, 1\times10^{-5},\, 1\times10^{-6}\}$
             บนกราฟเดียวกัน ที่ระยะคงที่ 10~cm
             (แกน $x$: เวลา (s), แกน $y$: ระยะทาง (cm))}
    \label{fig:exp3_scenarioA}
\end{figure}

\subsubsection{Scenario B: MSD Free Vibration --- Settling Time และ Variance}

Scenario B ทดสอบพฤติกรรมของ KF ในสภาวะ Dynamic โดยดึงมวลขึ้นจาก
ตำแหน่ง Equilibrium เป็นระยะ $x_0 \approx 7.2$~cm แล้วปล่อย (Free Vibration)
Settling Time นิยามเป็นเวลาที่ KF Output มี Amplitude ของการแกว่ง
ลดลงต่ำกว่า 5\% ของ Initial Amplitude และอยู่ในนั้นตลอดจนสิ้นสุดการบันทึก

\begin{table}[htbp]
    \centering
    \caption{ผลการวัด Settling Time และ Variance Reduction ของ KF
             ในสภาวะ MSD Free Vibration (Scenario B, 3 Trials)}
    \label{tab:exp3_scB_results}
    \begin{tabular}{clcccc}
        \hline
        \textbf{Trial}
            & \textbf{ค่า $Q$}
            & \textbf{Init.\ Amp.\ (cm)}
            & \textbf{Settle (ms)}
            & \textbf{Var KF (cm$^2$)}
            & \textbf{Noise ลดลง (\%)} \\
        \hline
        \multirow{3}{*}{1}
            & $1\times10^{-3}$ & 4.77 & 3232 & 1.2792 & 42.4 \\
            & $1\times10^{-5}$ & 4.59 & 3892 & 0.7948 & 64.2 \\
            & $1\times10^{-6}$ & 4.62 & 3892 & 0.7897 & 64.4 \\
        \hline
        \multirow{3}{*}{2}
            & $1\times10^{-3}$ & 4.95 & 3695 & 1.3223 & 40.3 \\
            & $1\times10^{-5}$ & 4.96 & 4223 & 0.7542 & 65.9 \\
            & $1\times10^{-6}$ & 5.07 & 4685 & 0.7717 & 65.1 \\
        \hline
        \multirow{3}{*}{3}
            & $1\times10^{-3}$ & 4.56 & 3894 & 1.3116 & 42.8 \\
            & $1\times10^{-5}$ & 4.32 & 5214 & 0.7182 & 68.7 \\
            & $1\times10^{-6}$ & 4.36 & 5280 & 0.7114 & 69.0 \\
        \hline
        \multirow{3}{*}{\textbf{Mean}}
            & $1\times10^{-3}$ & 4.76 & \textbf{3607} & 1.3044 & \textbf{41.8} \\
            & $1\times10^{-5}$ & 4.62 & \textbf{4443} & 0.7557 & \textbf{66.3} \\
            & $1\times10^{-6}$ & 4.68 & \textbf{4619} & 0.7576 & \textbf{66.2} \\
        \hline
    \end{tabular}
\end{table}

\begin{figure}[htbp]
    \centering
    % TODO: แทนด้วยกราฟจริงจาก MATLAB (plot ทั้ง 3 Q บนกราฟเดียวกัน)
    \includegraphics[width=0.85\textwidth]{exp3_scenarioB.png}
    \caption{Time-Domain Plot (Scenario B): Raw Data $z_k$ และ KF Output $\hat{x}_k$
             สำหรับ $Q$ ทั้ง 3 ค่า บนกราฟเดียวกัน
             ระบบ MSD Free Vibration เริ่มต้นที่ $x_0 \approx 7.2$~cm
             จาก Equilibrium
             (แกน $x$: เวลา (s), แกน $y$: ระยะทาง (cm))}
    \label{fig:exp3_scenarioB}
\end{figure}

\subsubsection{สรุปผล Quantitative}

\begin{table}[htbp]
    \centering
    \caption{สรุปค่า Quantitative เปรียบเทียบทั้ง 2 Scenario
             ($R = 0.0337$~cm$^2$, ค่าเฉลี่ยจาก 3 Trials)}
    \label{tab:exp3_summary}
    \begin{tabular}{lcccc}
        \hline
        \textbf{ค่า $Q$}
            & \textbf{$Q/R$}
            & \textbf{$K_{k,\mathrm{ss}}$}
            & \textbf{Settle Time (ms)}
            & \textbf{Var Reduction (\%)} \\
        \hline
        $1\times10^{-3}$ & $2.97\times10^{-2}$ & 0.136 & \textbf{3607} & 41.8 \\
        $1\times10^{-5}$ & $2.97\times10^{-4}$ & 0.017 & 4443          & \textbf{66.3} \\
        $1\times10^{-6}$ & $2.97\times10^{-5}$ & 0.005 & 4619          & 66.2 \\
        \hline
    \end{tabular}
\end{table}

\paragraph{หมายเหตุเกี่ยวกับหน่วย}
กราฟในรูปที่~6 และ~7 แสดงหน่วยแกน $y$ เป็น cm เพื่อความสะดวก
ในการอ่านค่าและเปรียบเทียบ ค่าดังกล่าวสามารถแปลงเป็นหน่วย SI (m)
ได้โดยคูณด้วยตัวประกอบ $1 \times 10^{-2}$
% \begin{figure}[htbp]
%     \centering
%     \includegraphics[width=0.9\textwidth]{figures/Figure_ErrorPlot.pdf}
%     \caption{รูปที่ 8: Error Plot แสดง Residual $e_k = z_k - \hat{x}_k$
%     สำหรับ $Q \in \{1\times10^{-3},\ 1\times10^{-5},\ 1\times10^{-6}\}$
%     ที่ระยะคงที่ 10~cm
%     (แกน $x$: เวลา (s), แกน $y$: Residual ($\times 10^{-2}$~m))}
%     \label{fig:error_plot}
% \end{figure}

%% ─────────────────────────────────────────────────────────────────────────────
\subsection{สรุปผลการทดลอง}

\begin{itemize}
    \item \textbf{Scenario A (Static):} KF ทุกค่า $Q$ สามารถลด Noise ได้อย่าง
          มีนัยสำคัญ (71--99\%) โดย $Q = 1\times10^{-6}$ ให้ Variance Reduction
          สูงที่สุดเฉลี่ย 90.2\% ขณะที่ $Q = 1\times10^{-3}$ ให้ 76.5\%
          สอดคล้องกับทฤษฎีที่ $K_{k,\mathrm{ss}}$ ต่ำกว่าจะกรอง Noise ได้มากกว่า
          แต่แลกกับการตอบสนองที่ช้าลง

    \item \textbf{Scenario B (Dynamic):} $Q = 1\times10^{-3}$ มี Settling Time
          เร็วที่สุดเฉลี่ย 3607~ms เนื่องจาก $K_{k,\mathrm{ss}} = 0.136$
          ทำให้ Filter ติดตาม Measurement ได้ไว แต่ Variance Reduction
          ต่ำที่สุดที่ 41.8\% ขณะที่ $Q = 1\times10^{-5}$ และ $Q = 1\times10^{-6}$
          ให้ Variance Reduction ใกล้เคียงกัน ($\approx 66\%$) แต่ $Q = 1\times10^{-5}$
          settle เร็วกว่า $Q = 1\times10^{-6}$ อย่างชัดเจน (4443 vs.\ 4619~ms)

    \item \textbf{ค่า $Q$ ที่แนะนำ:} $Q = 1\times10^{-5}$ ให้ผลสมดุลดีที่สุด
          ระหว่าง Smoothness (Var Reduction 66.3\%) และ Responsiveness
          (Settle Time 4443~ms) โดยไม่เสียสละ Noise Rejection มากเกินไป
          เหมือน $Q = 1\times10^{-3}$ และตอบสนองได้เร็วกว่า $Q = 1\times10^{-6}$
          อย่างมีนัยสำคัญในสภาวะ Dynamic สอดคล้องกับ Nominal Value
          ที่ Al Tahtawi (2018) เลือกใช้
\end{itemize}

%% ─────────────────────────────────────────────────────────────────────────────
\subsection{อภิปรายผล}

\begin{itemize}
    \item \textbf{กลไกที่ $Q$ ส่งผลต่อ Filter:}
          $Q$ บวกเข้าไปใน $P_k^{(-)}$ ทุกรอบผ่านสมการ~\eqref{eq:cov_predict_exp3}
          ทำให้ Kalman Gain $K_k = P_k^{(-)}/(P_k^{(-)} + R)$ สูงขึ้น
          และ State Update $\hat{x}_k = \hat{x}_k^{(-)} + K_k(z_k - \hat{x}_k^{(-)})$
          ดึง Output เข้าหา Measurement มากขึ้น
          ผลคือ Filter ตอบสนองเร็วขึ้นแต่ยอมรับ Noise จาก Sensor มากขึ้น
          ดังที่เห็นในตารางที่~\ref{tab:exp3_scB_results}
          ว่า $Q = 1\times10^{-3}$ มี Settle Time สั้นที่สุด
          แต่ Variance Reduction ต่ำที่สุด

    \item \textbf{$Q = 1\times10^{-5}$ vs.\ $Q = 1\times10^{-6}$:}
          ผลที่น่าสนใจคือทั้งสองค่านี้ให้ Variance Reduction ใกล้เคียงกันมาก
          ($66.3\%$ vs.\ $66.2\%$) แต่ Settle Time ต่างกัน 176~ms
          เหตุผลคือ $K_{k,\mathrm{ss}}$ ต่างกัน 3.4 เท่า (0.017 vs.\ 0.005)
          ใน Scenario Static ความต่างของ Variance นี้เล็กน้อยเพราะ Signal
          ไม่เปลี่ยนแปลง แต่ใน Dynamic Scenario ความต่างของ Responsiveness
          จะเห็นได้ชัดกว่า

    \item \textbf{$Q/R \ll 1$ ในทุกกรณี:}
          เนื่องจาก $R = 0.0337$~cm$^2$ ใหญ่กว่า $Q$ ทุกค่ามาก
          Filter จึงโน้มเอียงไปทาง ``เชื่อ Model'' เป็นค่า Default อยู่แล้ว
          ซึ่งแตกต่างจาก Al Tahtawi (2018) ที่ใช้ $R = 2.92\times10^{-3}$
          ทำให้ $Q = 1\times10^{-3}$ มี $Q/R > 1$ ในงานวิจัยนั้น
          แต่ \textbf{ทิศทางของ Effect เหมือนกันทุกประการ}
          คือ $Q$ ที่ใหญ่กว่าทำให้ Filter ตอบสนองเร็วกว่าเสมอ

    \item เมื่อเปรียบเทียบกับงานวิจัยของ Al~Tahtawi (2018) ซึ่งใช้
            $R = 2.92\times10^{-3}$~cm$^2$ พบว่าในงานวิจัยนั้น
            $Q = 1\times10^{-3}$ ทำให้ $Q/R > 1$ กล่าวคือ Filter เชื่อ
            Measurement มากกว่า Model ในขณะที่การทดลองนี้ใช้
            $R = 0.0337$~cm$^2$ ซึ่งใหญ่กว่า $Q$ ทุกค่า ทำให้ $Q/R < 1$
            ตลอด Filter จึงโน้มเอียงไปทาง Model มากกว่า อย่างไรก็ตาม
            ทิศทางของ Effect เหมือนกันทุกประการ กล่าวคือ $Q$ ที่ใหญ่กว่า
            ทำให้ Filter ตอบสนองเร็วกว่าเสมอ ไม่ว่า $Q/R$ จะมากกว่าหรือ
            น้อยกว่า 1 ความแตกต่างนี้เกิดจาก Sensor ที่ใช้ในการทดลองนี้
            มี Noise สูงกว่า (Variance ต่างกัน 11 เท่า) ดังนั้นการเปรียบเทียบ
            ค่าสัมบูรณ์ของ $Q$ จึงไม่มีความหมาย แต่ควรเปรียบเทียบจาก
            อัตราส่วน $Q/R$ แทน

    \item \textbf{สภาวะ Overconfident ($Q \to 0$):}
          หาก $Q = 0$ ค่า $P_k$ จะ Converge ลงสู่ศูนย์
          ส่งผลให้ $K_k \rightarrow 0$ ตามสมการ~\eqref{eq:scalar_gain_exp3}
          Filter จะหยุดรับข้อมูลใหม่จาก Sensor โดยสิ้นเชิง
          ซึ่งจะทำให้ KF ไม่สามารถติดตามการแกว่งของ MSD ได้เลย
          แม้แต่ $Q = 1\times10^{-6}$ ที่ใช้ในการทดลองนี้
          ก็ยังแสดงให้เห็นว่า Settle Time ช้ากว่า $Q = 1\times10^{-5}$
          อย่างมีนัยสำคัญในสภาวะ Dynamic

    \item \textbf{ข้อสังเกตเกี่ยวกับ Scenario B:}
          Settling Time ที่วัดได้ในการทดลองนี้สะท้อนทั้งการที่ KF ติดตาม
          Measurement ได้ไวแค่ไหน \textbf{และ} ระยะเวลาที่ระบบ MSD จริง
          Damp Out ทางกายภาพ ดังนั้นความต่างของ Settling Time ระหว่าง
          $Q = 1\times10^{-3}$ (3607~ms) กับ $Q = 1\times10^{-6}$ (4619~ms)
          บางส่วนมาจาก Filter Lag ที่ต่างกัน ไม่ใช่ทั้งหมดจาก Physical Damping
\end{itemize}

\subsection{ข้อเสนอแนะเพื่อการพัฒนา}
 
จากการวิเคราะห์ข้อจำกัดที่พบในการทดลอง มีข้อเสนอแนะเพื่อปรับปรุงและต่อยอดดังนี้
 
\begin{itemize}
    \item \textbf{ขยายช่วงค่า $Q$ ที่ทดสอบ:} การทดลองนี้ทดสอบเพียง 3 ค่าตาม Al Tahtawi (2018) ซึ่งอาจไม่เพียงพอในการหาจุดสมดุลที่เหมาะสมที่สุด ควรทดสอบค่า $Q$ เพิ่มเติมในช่วงระหว่าง $10^{-6}$ ถึง $10^{-3}$ เพื่อให้เห็น Trade-off ระหว่าง Smoothness และ Responsiveness ได้ละเอียดยิ่งขึ้น
    % \item \textbf{วัด Response Time อย่างเป็นระบบ:} การนิยาม Settling Time ผ่าน Envelope Method ในการทดลองนี้มีความซับซ้อน ควรออกแบบ Scenario Step Change ที่ชัดเจนกว่า เช่น เลื่อนวัตถุจาก 10 cm ไปยัง 20 cm อย่างรวดเร็ว เพื่อให้วัด Response Time แบบ 5\% Settling ได้ตรงไปตรงมา
    \item \textbf{พิจารณา Adaptive Process Noise:} ค่า $Q$ คงที่ตลอดเวลาทำให้ Filter ไม่สามารถปรับตัวเมื่อพลศาสตร์ของระบบเปลี่ยน เช่น เมื่อมวลเคลื่อนที่เร็วหรือช้า ในอนาคตอาจพิจารณาใช้ Adaptive Kalman Filter ที่ประมาณค่า $Q$ ใหม่ทุกรอบจาก Innovation Sequence แทน
\end{itemize}